% SPDX-License-Identifier: Apache-2.0
\section{The Channel as Hardware}
\label{sec:x-buffer}

Every lowered unit with a channel port has been checked against its
trace with the testbench playing the channel's other side. Between
two lowered units on a board there is no testbench, and there is no
channel either: each unit's port is one side of a handshake, and
wiring the sides together makes a combinational loop, since a
receiver that takes whatever is offered says \code{ready} when it
sees \code{valid}. The channel, the elastic buffer of two, must be a
piece of hardware of its own, and \code{Buffer} in
\code{txhdl\_parts} is that piece, written in the lowerable subset
with the runtime's rules and nothing else.

The example is its check. Each cycle the testbench sets the buffer's
three inputs, an offer and its word and a take, from a pattern that
fills the buffer, holds the sender back, drains it, and offers and
takes in the same cycle both empty and full; it does the same offer
and take on a runtime channel; and after the cycle it asserts that the
buffer's \code{ready}, \code{valid} and head are the channel's. The
buffer's state is what is compared, as the edge leaves it, since a
unit's output wires are set in the step from the state before it and
show the same a step later, which is the netlist check's business:
the netlist is checked against this trace under nvc and Verilator as
every lowered unit is. So the two checks together say that a channel
between two lowered units on a board does what the channel between
the two units in the run did.

\begin{figure}[htbp]
\centering
\input{buffer_timing.tex}
\caption{The buffer: offers while there is room, takes every other
cycle, then both; \code{head\_full} and \code{tail\_full} are its
two slots.}
\label{fig:x-buffer}
\end{figure}

\lstinputlisting[caption={\code{ex\_buffer.rs}. Run with \code{bazel run //lib/examples:ex\_buffer}.},label={lst:x-buffer}]{lib/examples/ex_buffer.rs}

\lstinputlisting[style=ascii,caption={What \code{ex\_buffer} printed when the build ran it: the pattern and the buffer's state after each cycle, then the lowering.},label={lst:x-buffer-out}]{out_buffer.txt}
